\documentclass[class=article,border=10pt, margin=3mm]{standalone}
\usepackage{verbatim}
\usepackage{textcomp}
\usepackage{tikz}
\usetikzlibrary{shapes,arrows}
\usepackage[americanresistors,americaninductors,RPvoltages]{circuitikz}
\usepackage{amsmath}


\begin{document}

\begin{circuitikz}[american, scale=1.2]
    
    % 1. Definición de Coordenadas del Puente (Ajusta el 1.5 para cambiar la altura)
    \coordinate (T) at (3, 1.5);  % Nodo Superior
    \coordinate (B) at (3, -1.5); % Nodo Inferior
    \coordinate (L) at (1.5, 0);  % Nodo Izquierdo
    \coordinate (R) at (4.5, 0);  % Nodo Derecho


    \draw (0, 1.5) to[vsourcesin, l=$V_{gen}$] (0, -1.5);
    \draw (0, 1.5) -- (T);   
    \draw (0, -1.5) -- (B); 

   
    \draw (L) to[D*] (T); 
    \draw (L) to[D*] (B); 
    \draw (T) to[D*] (R); 
    \draw (B) to[D*] (R); 

   
    \draw (L) -- ++(-0.3, 0) node[ground] {};

   
    \draw (R) -- (5.5, 0) 
        to[C, l_=$100\,\mu\text{F}$, *-] (5.5, -2.5) 
        node[ground] {};

    
    \draw (5.5, 0) -- (7.0, 0)
        to[R, name=res, *-] (7.0, -2.5) 
        node[ground] {};

    \node[anchor=west] at (6.1, -0.75) {$R_L$};
    \node[anchor=east] at (6.8, -1.25) {$4.7\,\text{k}\Omega$};
   

    \node[anchor=west] at (7.4, -0.5) {$+$};
    \node[anchor=west] at (7.4, -1.25) {$V_{out}$};
    \node[anchor=west] at (7.4, -2.0) {$-$};

\end{circuitikz}

\end{document}